\chapter{3D打印与快速原型}
\label{chap:additive-manufacturing}
\label{ch:additive}

\nt{本章复习的核心是\textbf{材料以什么形式送入、如何在一层内选区成形、层间怎样结合、何处需要支撑、成形后还要做什么}。名称中的“激光”“喷射”不能单独决定工艺类型，应沿材料与能量的传递过程判断。}

\section{增材制造的思路与数据流程}
\dfn{3D打印与快速原型}{3D打印依据数字模型，通过逐层添加并结合材料形成三维物体。快速原型（Rapid Prototyping，RP）强调快速制作实物原型的用途；增材制造（Additive Manufacturing）强调逐层增加材料的制造方式。原型可用于外观、功能、装配和设计交流，部分工艺也可直接制造功能零件。}

传统减材制造从毛坯去除材料，增材制造逐层构建所需形状，适合复杂几何、个性化和小批量等任务。它通常减少专用模具需求，但仍需考虑基板、定位、支撑、热变形和后加工；不能把讲义中的“无需工装”理解为任何工艺都完全没有装夹或辅助结构。

基本流程为
\[
\text{CAD模型}\to\text{数据转换与修复}\to\text{摆放和支撑}
\to\text{切片与路径}\to\text{打印}\to\text{取件及后处理}.
\]
打印前需确认单位、尺寸和成形方向，设置层厚、材料与工艺参数。打印后可能包含去支撑、清粉、清洗、固化、烧结、热处理或机加工，\textbf{打印完成不等于零件已经满足使用要求}。
\sourcepage{第7章}{3--7}

\section{STL：用三角面片传递形状}
STL以三角形逼近模型表面，每个面片记录三个顶点和一个面法向量，既有ASCII文本形式，也有二进制形式。它表达离散几何，不保留原CAD中的参数化特征、精确解析曲面或完整设计约束。传统STL还不提供可靠的单位语义，导入导出时应核对尺寸。

\subsection{面片精度与可打印性}
曲面离散会引入几何近似误差。三角片过粗，圆弧和曲面会产生明显折线；片过密，数据增大而不一定提高设备能达到的实际精度。应使离散误差与所需尺寸、表面质量及制造能力相匹配。

\begin{itemize}
\item \textbf{间隙与孔洞：}边界未闭合，切片软件难以判定内部区域。
\item \textbf{面法向翻转：}局部内外方向不一致；顶点绕序与法向也应一致。
\item \textbf{重叠、自相交与非流形连接：}可能造成截线分支、重复区域或无法确定实体内部。
\item \textbf{薄壁、小孔与封闭空腔：}即使网格拓扑正确，也可能超过制造能力，或无法去除支撑与残留材料。
\end{itemize}
判断STL质量时，应同时检查\textbf{几何有效性}与\textbf{制造可行性}。水密模型可以用于切片，但水密本身不能保证薄壁会成形、孔会保持通畅。
\sourcepage{第7章}{8--10}

\section{层厚、成形方向、支撑与路径规划}
\subsection{阶梯效应：降低层厚并非没有代价}
分层制造用一层层截面逼近连续曲面。倾斜或弯曲表面在层与层之间出现台阶，称为阶梯效应；其程度与层厚、局部表面倾角及成形方向有关。减小层厚通常能减轻该效应，却增加层数和逐层操作时间。
\reviewfig[.85]{ch07-staircase.pdf}{不同层厚下的阶梯近似。摘自第7章讲义，PDF第11页。}{fig:printing-staircase}

\ex{层厚与层数的估算}{某零件沿成形方向的高度为 $H=60\,\mathrm{mm}$，采用均匀层厚 $h$。忽略首层和末层调整，层数约为
\[
N\approx\left\lceil\frac{H}{h}\right\rceil.
\]
$h=0.2\,\mathrm{mm}$ 时约300层，改为 $0.1\,\mathrm{mm}$ 时约600层。若每层固定准备动作耗时 $t_0$，这部分时间由 $300t_0$ 增为 $600t_0$；\textbf{总打印时间不能直接认定为恰好翻倍}，因为扫描、挤出、曝光、移动和冷却时间还由截面与参数决定。}
\sourcepage{第7章}{11、13--17}

\subsection{支撑为何必要，为什么有些粉末工艺不需独立支撑}
支撑可承担悬垂区域、保持与基板的连接、限制变形；在某些金属工艺中还承担散热与锚固作用。图\ref{fig:printing-support}区分以周围牺牲材料包围形状的互补支撑和另外生成的显式支撑。
\reviewfig[.88]{ch07-support.pdf}{互补支撑与显式支撑的示意。摘自第7章讲义，PDF第12页。}{fig:printing-support}

松散粉末可以托住正在成形的区域，因此典型聚合物SLS通常不需要另加几何支撑；但仍需考虑热变形与粉末清理。\textbf{不能推广为所有粉末床工艺均不需要支撑}，SLM金属零件常仍需连接基板的支撑。

支撑越多，耗材、打印时间、去除时间和接触处表面损伤通常越大。设计时需平衡可靠性与后处理成本，并保留去支撑或排粉通道。讲义中的45$^\circ$可作特定情境的经验提示，\textbf{允许悬垂角不是所有材料与设备通用的定值}。
\sourcepage{第7章}{12、16、19、24--25}

\subsection{工艺规划是多目标问题}
成形方向同时影响高度、层数、支撑面积、关键面台阶和力学各向异性。降低高度可能缩短逐层准备时间，却不一定减少支撑；将重要受力方向横跨薄弱层间结合方向，也可能降低承载能力。

切片后还需规划轮廓与内部填充的扫描或挤出路径，选择方向、间距、顺序、起止点与空行程，并设置温度、速度、功率等工艺参数。路径关系到成形时间、表面质量、热积累和强度。因而不能只用单一“打印速度”评价工艺效率。
\sourcepage{第7章}{13--17}

\section{光聚合物成形：SLA、DLP与PolyJet}
\subsection{SLA与DLP：同为光固化，曝光方式不同}
\dfn{SLA}{光固化成形（Stereolithography，SLA）以液态光敏树脂为材料，使激光按切片截面选择性扫描并固化，再进行下一层的成形。局部固化区域逐步构成三维零件。}
\reviewfig[.82]{ch07-sla.pdf}{SLA的光源、扫描与树脂成形区。摘自第7章讲义，PDF第18页。}{fig:printing-sla}

\textbf{DLP（Digital Light Processing）}利用投影图案使一层所需区域曝光，区别于典型SLA沿路径逐点扫描。整层投影可减少逐点扫描的时间，但单层固化后仍有铺层、抬升、剥离等动作；实际效率还取决于曝光时间、截面、设备及材料。

两者适于表达较细的形状和较好的表面细节，但材料性能、光固化收缩、支撑接触和后固化会影响尺寸及性能。典型后处理包括取件、清洗、去支撑和继续固化。\textbf{分辨率、尺寸准确度和表面粗糙度是不同指标}，不能因为像素或光斑较小就认定零件所有尺寸都同样准确。
\sourcepage{第7章}{18--20、48}

\subsection{PolyJet：喷射材料，再以光固化}
PolyJet先由喷头将液态光聚合物微滴喷到成形区，再以紫外光固化，薄层重复堆积。多喷头可组合不同材料、颜色或支撑材料，适合多材料与外观原型。支撑材料与本体材料可同时沉积，但成形后仍需去除支撑。

它与DLP都可使用光聚合物，但\textbf{材料供给方式不同}：DLP在树脂槽中选择性曝光；PolyJet先选择性喷射材料，再固化。也不要与3DP混淆，后者是向已有粉末床喷粘结剂。
\sourcepage{第7章}{21}

\section{粉末床激光成形：SLS与SLM}
\subsection{SLS：选区使粉末结合成层}
选择性激光烧结（Selective Laser Sintering，SLS）将粉末铺成薄层，经预热后，激光按截面选区扫描，使材料结合并与前一层连接；随后下降工作台、重新铺粉，反复完成零件。具体结合机制取决于材料与工艺参数，不应把所有材料都理解成相同的“颗粒半熔化”过程。

典型聚合物SLS利用周围未成形粉末承托复杂形状，不需另设传统显式支撑。它便于形成复杂内部结构，但表面会受粉末粒度与成形机制影响；孔隙、热收缩和各向异性也需要考虑。后处理包括冷却、取出粉床、清除松散粉末及必要的表面处理，细长孔、盲孔和封闭内腔尤其难清粉。

\subsection{SLM：金属粉末的选区熔化}
选择性激光熔化（Selective Laser Melting，SLM）用激光使指定区域的金属粉末熔化，凝固后形成冶金结合，再逐层铺粉成形。它可用于金属功能件和复杂模具，但需考虑热梯度、残余应力、翘曲、支撑和后续处理。

\begin{table}[H]
\centering\small
\begin{tabularx}{\linewidth}{@{}p{2cm} X X@{}}
\toprule
比较项目 & 典型聚合物SLS & SLM金属粉末成形 \\
\midrule
材料与结合 & 常用聚合物粉末，激光使局部粉末结合。 & 金属粉末选区熔化并凝固。 \\
支撑需求 & 周围粉末通常提供几何承托。 & 粉末承托不能代替锚固、导热和变形控制所需支撑。 \\
后处理重点 & 冷却、清粉、表面处理等。 & 去支撑、与基板分离，按要求热处理、表面与精密加工等。 \\
性能判断 & 取决于材料、致密程度、成形方向及参数。 & 可得到高致密金属件，但不能保证任意参数下都无孔隙、无缺陷。 \\
\bottomrule
\end{tabularx}
\end{table}

\nt{SLM与SLA都按切片逐层选区成形，但前者\textbf{熔化金属粉末}，后者\textbf{光固化液态树脂}。二者不是同一种材料反应；讲义中的类比应理解为逐层制造流程上的相似。比较SLS与SLM时也不能只凭名称判断所有零件的强度高低。}
\sourcepage{第7章}{22--26}

\section{薄片与挤出成形：LOM、FDM}
\subsection{LOM：层间粘接与轮廓切割}
叠层实体制造（Laminated Object Manufacturing，LOM）把薄片或箔材逐层粘接，再按切片轮廓切割；纸张、塑料或某些复合薄片可作材料。典型装置中，热压辊使新薄片与已成形部分结合，激光或刀具切割轮廓，多余材料在制造中可保留在原位提供支撑。

其特点是薄片供料、分层粘接和轮廓切割，适于某些较大尺寸的形状模型；代价是余料去除，内部封闭结构的余料可能无法清除。表面与尺寸稳定性也受薄片和粘接层影响。\textbf{“周围余料能支撑”并不意味着“任意内腔都能完成后处理”。}
\sourcepage{第7章}{27--29}

\subsection{FDM：挤出、铺丝与层间结合}
熔融沉积成形（Fused Deposition Modeling，FDM）将聚合物丝材送入加热区，材料软化或熔融后经喷嘴挤出，按路径铺放并与已有材料结合，冷却后保持形状。牵引滚轮提供送料，喷嘴的相对运动决定铺丝路径；直角坐标与Delta等机构只是实现运动的不同架构。
\reviewfig[.35]{ch07-fdm.pdf}{FDM的送料、加热与铺丝过程。摘自第7章讲义，PDF第30页。}{fig:printing-fdm}

常见优点是供料与设备结构直观、材料选择较广，适于原型及适当的功能件。表面可见层纹，尺寸受挤出宽度、冷却收缩和工艺参数影响；层间结合与丝路方向使力学性能具有方向性。多喷头可以实现多色、多材料或可溶支撑，但需处理喷头切换、材料相容性和支撑去除。

\nt{不能把“材料种类多”“成本低”或“强度好”理解为所有FDM零件都优于其他聚合物工艺。应先规定材料、几何、成形方向、测试方向和质量要求，再作比较。}
\sourcepage{第7章}{30--32}

\section{粘结剂喷射与增减材复合：3DP、SDM}
\subsection{3DP：在粉末床上喷粘结剂}
本章讲义中的狭义3DP指：铺设粉末后，由喷头向所需截面喷射液态粘结剂，使该处粉末结合；再下降工作台、铺新粉和喷射，循环形成零件。粉末可以是石膏、陶瓷或金属等，具体后处理与可获得性能随材料体系变化。

未粘结粉末可在制造过程中承托几何，但打印后的初始坯体可能较脆弱。取件与清粉后，根据目标可能还需浸渗、脱脂或烧结等处理，不能笼统认定“刷除粉末就获得最终功能件”。

\textbf{区分三种“喷射”：}PolyJet喷的是构成物体的光聚合物；3DP喷的是把粉末结合起来的粘结剂；FDM是经喷嘴连续挤出热塑性材料，不是向粉床喷粘结剂。
\sourcepage{第7章}{33--34、49}

\subsection{SDM：沉积与切削交替完成形状}
形状沉积制造（Shape Deposition Manufacturing，SDM）交替进行材料沉积、数控加工及必要的应力处理，再继续沉积。支撑材料与本体材料按需要构建，最后去除支撑。其思路是用增材形成体积，再用减材改善局部形状和表面质量。
\reviewfig[.9]{ch07-sdm.pdf}{SDM将材料沉积与切削结合。摘自第7章讲义，PDF第35页。}{fig:printing-sdm}

复合工艺仍需考虑刀具可达性、装夹基准与多道工序间的误差。某些表面在被后续材料覆盖之前加工，可能比整个零件完成后再进入刀具更容易；工序顺序因此直接影响可制造性。
\sourcepage{第7章}{35}

\section{工艺选择与典型应用}
\begin{table}[H]
\centering\small
\begin{tabularx}{\linewidth}{@{}p{2.2cm} X X@{}}
\toprule
工艺 & 识别关键词 & 复习时首先检查 \\
\midrule
SLA / DLP & 树脂槽；扫描光点 / 层图案投影 & 树脂性能、支撑、清洗和后固化。 \\
PolyJet & 光聚合物喷射后固化 & 多材料能力、支撑与材料相容性。 \\
SLS / SLM & 粉末床激光结合 / 金属熔化 & 清粉、热变形；金属支撑与后处理。 \\
LOM & 薄片粘接并切轮廓 & 层间结合、余料可去除性。 \\
FDM & 热塑性丝材熔融挤出 & 丝路、层间结合、收缩与支撑。 \\
3DP & 粉末床上喷粘结剂 & 坯体强度、清粉及后续致密化。 \\
SDM & 沉积和切削交替 & 工序次序、刀具可达性与基准。 \\
\bottomrule
\end{tabularx}
\end{table}

\ex{有内腔的复杂零件如何比较工艺}{若目标是演示外观与装配关系，可重点比较形状细节、成本、颜色、尺寸与支撑痕迹；若目标是金属承载件，则须进一步考察材料、缺陷、残余应力、后处理和质量验证。无论哪种工艺，都要检查内腔中的\textbf{支撑、树脂、余料或粉末能否排出}。把内腔在CAD中画出来，并不代表它一定能打印并清理完成。}

典型应用包括外观与功能原型、艺术复制、个性化产品、复杂模具及轻量化构件。讲义中的医学模型展示了“断层数据 $\to$ 分割与重建 $\to$ CAD设计 $\to$ 实体模型”的联系：可视化模型、手术规划模型和最终植入物的用途不同，所需材料与验证要求也不同，不能仅凭几何形状相同就互相替代。
\sourcepage{第7章}{36--47}

\Needspace{11cm}
\section{复习自测与易错点}
\begin{enumerate}
\item 减小层厚为什么通常更慢？\textit{答：在相同成形高度下层数增加，逐层准备、曝光或扫描及运动等动作增多。}
\item STL面片更密，是否一定使实物更准确？\textit{答：否。离散误差只是总误差的一部分，还受设备、材料与工艺参数限制。}
\item SLA和DLP的主要区别是什么？\textit{答：典型SLA按路径扫描曝光，DLP按一层图案投影曝光。}
\item PolyJet和3DP的喷射材料有何不同？\textit{答：前者喷光聚合物等成形材料并固化，后者在粉末床喷粘结剂。}
\item SLS通常不需独立支撑，SLM是否也一样？\textit{答：不能据此推断。金属SLM常需支撑承担锚固、散热和变形控制。}
\item 打印方向只影响台阶吗？\textit{答：还影响层数、支撑、成形时间、关键尺寸和力学性能方向性等。}
\item LOM中周围余料的好处和代价是什么？\textit{答：制造时可支撑形状，完成后必须去除，复杂内腔可能难以清理。}
\item SDM为何要在沉积过程中穿插切削？\textit{答：可改善尺寸和表面，并利用尚未被覆盖时的刀具可达性。}
\end{enumerate}
\nt{复习比较题时，按“\textbf{材料形态 $\to$ 选区成形机制 $\to$ 层间结合 $\to$ 支撑 $\to$ 后处理 $\to$ 适用目标}”作答。避免只背工艺名称或绝对化的“最快、最便宜、最强”；条件明确的比较才有意义。}
